\documentclass[11pt, a4paper]{article}      
\usepackage[T2A]{fontenc}
\usepackage[utf8x]{inputenc}
\usepackage[russian]{babel}
\usepackage{csquotes}
\usepackage{amsmath}
\usepackage{amsthm}
\usepackage{enumitem}
\usepackage{amssymb}
\usepackage{listings}
\usepackage[pdftex]{graphicx}
\usepackage{caption}
\usepackage[a4paper,left=2cm,right=2cm,
top=2cm,bottom=2cm,bindingoffset=0cm]{geometry}
\usepackage[colorlinks=true, linkcolor=black, urlcolor=blue, citecolor=black]{hyperref}
\usepackage{indentfirst}
\usepackage[nottoc]{tocbibind}
\usepackage{arydshln}
\usepackage{listings}
\usepackage{color}
\usepackage{minted}

\definecolor{codegreen}{rgb}{0,0.6,0}
\definecolor{codeorange}{rgb}{0.8, 0.5, 0}

\captionsetup[figure]{name=Рисунок}

\def\newsection#1{
    \section{#1}
    % \addcontentsline{toc}{section}{#1}
}
\def\newsubsection#1{
    \subsection{#1}
    % \addcontentsline{toc}{subsection}{#1}
}

\lstdefinestyle{customc}{
  inputencoding=utf8x,
  extendedchars=\true,
  belowcaptionskip=1\baselineskip,
  breaklines=true,
  frame=L,
  xleftmargin=\parindent,
  language=C++,
  showstringspaces=false,
  basicstyle=\footnotesize\ttfamily,
  keywordstyle=\bfseries\color{blue},
  commentstyle=\color{codegreen},
  identifierstyle=\color{black},
  stringstyle=\color{codeorange},
  keepspaces=true,
}
\lstset{style=customc}

\begin{document}

\begin{center}
Министерство науки и высшего образования Российской Федерации \\
Федеральное государственное автономное образовательное учреждение \\
высшего образования\\ 
\textbf{«Национальный исследовательский \\ Нижегородский государственный университет \\
им. Н.И. Лобачевского (ННГУ)»}\\[1.5cm]
\textbf{Институт информационных технологий, математики и механики }\\[5.5cm]

\textbf{Отчет по лабораторной работе} \\
\textbf{по курсу}  \\
``Параллельные численные методы'' \\
``Блочное разложение Холецкого'' \\ [5cm]


\begin{flushright}
\begin{minipage}{0.52\textwidth}
\begin{flushleft}
\textbf{Выполнил:} \\
студент группы 3825М1ПМкн1 \\
Сучков В.Н.\\[0.3cm]
\textbf{Проверил:}\\
д.т.н., проф. кафедры МОСТ\\  Баркалов К.А.\\[0.3cm]
\end{flushleft}
\end{minipage}
\end{flushright}

\vfill

Нижний Новгород \\
2025

\thispagestyle{empty}
\end{center}

\newpage
\tableofcontents
\thispagestyle{empty}
\newpage
\newsection{Введение}
\setcounter{page}{1}
В рамках данной лабораторной работы был реализован и ислледован параллельный блочный алгоритм разложения Холецкого для
симметричной положитеьльно определённой матрицы. Особенностью данного алгоритма явялется то, что, помимо того, что они явялется
параллельным, в нём применяется разбиение входных данных на блоки, для повышения производительности за счёт более эффективного
использования кэш-памяти процессора.


Для параллелизации алгоритма использовался фреймворк OpenMP, а также в реализации алгоритма применялись SIMD интринсики из набора инструкций AVX2.


\newsection{Постановка задачи}
На языке C++ реализуется функция, выполняющая блочное разложение Холецкого вида 
\[
    A = LL^T,
\]
где $L$-нижнетреугольная матрица и $A = A^T > 0$.

\begin{lstlisting}[language=C++, frame=single]
    void Cholesky_Decomposition(double* A, double* L, int n);
\end{lstlisting}
где в свою очередь
\begin{itemize}
    \item $n$ - порядок матрицы $A$.
    \item $A$ - указатель на симметричную положитеьльно определённую матрицу размера $n \times n$ (row-major формат хранения).
    \item $L$ - указатель на результирующую нижнетреугольную матрицу (так же row-major формат хранения).
\end{itemize}

\textbf{Критерии и ограничения:}
\begin{itemize}
    \item Лимит по времени 100 с, по памяти: 128 Мб.
    \item Корректность ответа проверяется с помощью относительной невязки. Решение считается удовлетворительным если верно $\frac{|| LL^T - A ||_2}{||A||_2} < 0.01$.
    \item Показатель эффективности параллельной версии должен быть не менее 50\% для всех тестов.
\end{itemize}

\newsection{Описание метода решения задачи}
Основой для решения задачи послужит блочный алгоритм разложения Холецкого. Входная матрицы $A$ и $L$ разбиваются на блоки следующим образом
\[
    A =
    \left(\begin{array}{c c}
        A_{11} & A_{12} \\
        A_{21} & A_{22} \\
    \end{array}\right),
    L =
    \left(\begin{array}{c c}
        L_{11} & 0 \\
        L_{21} & L_{22} \\
    \end{array}\right),
\]
Блок $A_{11}$ имеет размер $k \times k$, где $k$ - некоторый выбранный заранее размер блока. Тогда, матрицу $A$ можно 
записать как произведение матриц $LL^T$, чтобы найти неизвестный блоки матрицы $L$.
\[
    A =
    \left(\begin{array}{c c}
        A_{11} & A_{12} \\
        A_{21} & A_{22} \\
    \end{array}\right) =
    \left(\begin{array}{c c}
        L_{11} & 0 \\
        L_{21} & L_{22} \\
    \end{array}\right)
    \left(\begin{array}{c c}
        L_{11}^T & L_{21}^T \\
        0 & L_{22}^T \\
    \end{array}\right) =
    \left(\begin{array}{c c}
        L_{11}L_{11}^T & L_{11}L_{21}^T \\
        L_{21}L_{11}^T & L_{21}L_{21}^T + L_{22}L_{22}^T \\
    \end{array}\right)
\]
Отсюда можем записать несколько равенств:
\begin{itemize}
    \item $L_{11}L_{11}^T = A_{11}$ - из этого равенства зная $A_{11}$ найдём $L_{11}$ с помощью неблочного разложения Холецкого.
    \item $L_{21}L_{11}^T = A_{21}$ - получив $L_{11}^T$ из равенства выше сможем найти $L_{21}$ как решение набора СЛАУ.
    \item $L_{21}L_{21}^T + L_{22}L_{22}^T = A_{22}$ - в этом уравнении сделаем замену вида $\tilde{A} = A_{22} - L_{21}L_{21}^T$. Тогда уравнение
    запишется в виде $\tilde{A} = L_{22}L_{22}^T$. Теперь, нахождение $L_{22}$ сводится к выполнению блочного разложения Холецкого к известной
    симметричной матрице $\tilde{A}$. 
\end{itemize}
Таким образом, получаем рекурсивный алгоритм, который последовательно обрабатывает матрицу $A$.


Суммарная трудоемкость алгоритма составляет $\frac{1}{3}n^3 + O(n^2)$ арифмитических операций. Метод Холецкого обладает
довольно высокой вычислительной устойчивостью. Рассмотрим решение с некоторой машинной погрешностью $(A + \Delta A)$. Норма матрицы
возмущений имеет верхнюю оценку
\[
    \| \Delta A \| \leqslant 3n^2\varepsilon_m \|A\| + O(\varepsilon_m^2),
\]
где $\varepsilon_m$ - машинная погрешность.

Сами матрицы хранятся в оперативной памяти компьютера в линейном представлении, как непрерывные последовательности чисел,
разложенные по строкам (row-major формат).

\newsection{Описание реализации алгоритма}
\newsubsection{Структуры данных}
Входная матрица $A$ и результирующая $L$ хранятся в памяти как одномерные массивы типа \texttt{double}.


Внутри программы используются ещё некоторый вспомогательные структуры: \texttt{MatBuffer} и \texttt{SubMatrix}, которые
используются для создания выровненного буфера памяти и хранения информации о матрицах и их подматрицах без лишнего копирования.
\begin{itemize}
    \item \texttt{\textcolor{blue}{class} MatBuffer} используется для аллокации и деаллокации выровненного буфера, который нужен для размещения в памяти
        промежуточных матриц, а также для выделения памяти в рамках тестирования.
    \item \texttt{\textcolor{blue}{struct} SubMatrix} хранит в себе указатель на начало матрицы в памяти, её размеры, а также расстояние в памяти между соседними строками,
        что позволяет осуществлять доступ к отдельным подматрицам некоторой матрицы путем задания начала подматрицы, её размеров и шага между соседними строками.
\end{itemize}


\newsubsection{Описание логики программы}
Работу функции можно разбить на следующие этапы:
\begin{itemize}
    \item [\textbf{1.}] Разложение верхнего левого блока $A_{11}$ размера $k \times k$ с помощью параллельного неблочного метода Холецкого.
    \item [\textbf{2.}] Параллельное решение СЛАУ с $A_{21}$ и $L_{11}$ для нахождения $L_{21}$.
    \item [\textbf{3.}] Параллельное матричное умножение $L_{21}L_{21}^T$ и нахождение матричной разности $\tilde{A} = A_{22} - L_{21}L_{21}^T$.
    \item [\textbf{4.}] Реккурсивный вызов функции для разложения $\tilde{A}$.
\end{itemize}

\textbf{1. Разложение верхнего диагонального блока.} Для разложения верхнего диагонального блока используется
параллельный алгоритм Холецкого. В неблочном алгоритме сначала вычисляется диагональный элемент матрицы $L$, после
чего параллельно вычисляются элементы под диагональным. Параллельность достигается за счёт использования директивы
\texttt{\#pragma omp parallel for}. При этом для вычисления сумм элементов матрицы $L$ используются
в том числе и SIMD инструкции. За счёт использования FMA (Fused-multiplication-addition) инструкций получается утилизировать параллелизм на уровне данных и 
обрабатывать числа пачками по 4 за одну итерацию внутреннего цикла параллельной секции.

\textbf{2. Параллельное решение СЛАУ.}  Для решения набора СЛАУ используетя алгоритм Гаусса. Параллельность тут достигается за 
счёт решения нескольких систем одновременно (для разных правых частей). Также, здесь используются SIMD FMA инструкции.

\textbf{3. Параллельное матричное умножение.} Матричное умножение выполняется по классическому алгоритму, однако, в данном
случае можно заметить, что умножение идёт матрицы на транспонированную матрицу, что означает доступ к обоеим матрицам вдоль строк,
что сокращает число кэш промахов. В данном случае снова для нахождения сумм произведений используются SIMD FMA инструкций, а так же
параллелизуется проход по разным строкам.

\newsection{Реализация параллельного алгоритма с импользованием общей памяти}
В неблочном алгоритме Холецкого основной блок вычислений не может быть параллельным (ввиду зависимости между итерациями цикла),
вместо этого параллельно выполняются вычисления поддиагональных элементов
\begin{lstlisting}[language=C++, frame=single]
    for (int j = 0; j < A.m; ++j) {
        // Находим диагональный элемент (j, j) в L.

#pragma omp parallel for private(...) schedule(dynamic, 8)
        for (int s = j + 1; s < L.n; ++s) {
            // Находим поддиагональные элементы используя найденный выше элемент
            // даигонали.
        }
    }
\end{lstlisting}

Для вычисления решения СЛАУ основной блок помещается под директиву \texttt{\#pragma omp parallel for}
\begin{lstlisting}[language=C++, frame=single]
#pragma omp parallel for
    for (int k = 0; k < B.n; ++k) {
        for (int j = 0; j < A.m; ++j) {
            for (int s = 0; s < quot; ++s) {
                // Непосредственно вычисления значений выходной марицы.
        }
    }
}
\end{lstlisting}

Матричное произведение вычисляется схожим образом с решением СЛАУ.

В итоговом блочном алгоритме последовательно выполняется набор параллельных алгоритмов, как было описано в начале раздела.

\newsection{Описание тестовых задач}
В рамках изучения реализованного алгоритма была проведена проверка корректности результата вычислений и замеры производительности.

\newsubsection{Проверка корректности результата}
Для проверки корректности используется относительная невязка $\frac{|| LL^T - A ||_2}{||A||_2}$, а вернее её оценка с использованием
матричной нормы Фробениуса $\frac{|| LL^T - A ||_2}{||A||_2} \leqslant \frac{\sqrt{n}|| LL^T - A ||_F}{||A||_F}$, где $n$ - порядок матрицы $A$.
В рамках тестов размер блока в блочном алгоритме был выбран $k = 64$. Такой размер позволяет эффективно использовать кэш память процессора
и уменьшить число промахов в кэш.

Для тестов корректности и производительности специально была написана функция для генерации случайных входных данных. Чтобы проверить, что 
алгоримт корректно отрабаывает, была произведена серия тестов с вычислением относительной погрешности.

\begin{table}[H]
    \centering
    \begin{tabular}{| c | c |}
        \hline
        Порядок входной матрицы $n$ & Относительная погрешность \\
        \hline
        4 & $8.324 \times 10^{-17}$ \\
        7 & $1.351 \times 10^{-17}$ \\
        13 & $2.383 \times 10^{-16}$ \\
        31 & $3.131 \times 10^{-16}$ \\
        54 & $3.953 \times 10^{-16}$ \\
        119 & $8.698 \times 10^{-16}$ \\
        128 & $9.524 \times 10^{-16}$ \\
        137 & $1.013 \times 10^{-15}$ \\
        153 & $1.125 \times 10^{-15}$ \\
        201 & $1.477 \times 10^{-15}$ \\
        997 & $6.312 \times 10^{-15}$ \\
        \hline
    \end{tabular}
    \caption{Результаты проверки на корректность}
    \label{table:1}
\end{table}

Как можно видеть из таблицы \ref{table:1}, на всех испробованых размерах входных данных величина относительной невязки не превосходила
установленного критерия $< 0.01$.

\newsubsection{Тесты производительности и масштабируемости}
Для замеров производительности и эффективности была выбрана матрица с фиксированным размером $n = 2048$ и размером блока $k = 512$. Большой размер удобен тем,
что позволяет создать нагрузку на потоки и сократить относительное влияние накладных расходов на параллелизм. Оцениваться будут параметры ускорения
($S = \frac{T_1}{T_p}$) и эффективности ($E = \frac{S}{p} \times 100\%$), где $p$-число используемых потоков.

\begin{table}[H]
    \centering
    \begin{tabular}{| c | c | c | c |}
        \hline
        Число потоков $p$ & Время выполнения, мс & Ускорение $S$ & Эффективность $E$, \% \\
        \hline
        1 & 444.3 & 1 & 100 \\
        2 & 300.1 & 1.48 & 74.02 \\
        3 & 227 & 1.96 & 65.24 \\
        4 & 184 & 2.415 & 60.36 \\
        6 & 133.5 & 3.32 & 55.47 \\
        \hline
    \end{tabular}
    \caption{Результаты проверки производительности}
    \label{table:2}
\end{table}

Из таблицы можно сделать вывод, что при использовании небольшого количества потоков (2-4) наблюдается существенный рост производительности.
Однако с дальнейшим ростом числа потоков эффективность начинает падать. Во многом такая деградация ускорения вызвана исчерпанием
пропускной способности памяти. Во всех рассмотренных случаях, тем не менее, эффективность была выше порога в 50\%.

\newsection{Заключение}
В ходе данной лабораторной работы, был реализован блочный алгоритм Холецкого. Алгоримт был исследован на корректность, производительность
и масштабируемость. В ходе тестирования было показано, что алгоритм соотвествует требованиям на корректность с большим запасом (значения
относительно невязки порядка $10^{-16}$). С точки зрения производительности алгоритм соотвествует ожиданиям, показывая 
неплохую масштабируемость ($>50$\% эффективности на всех многопоточных тестах) и высокую скорость работы.

\begin{thebibliography}{99}

\bibitem{Barkalov}
    Баркалов К.А. Метод Холецкого: презентация лекции по курсу ``Параллельные численные методы''. ННГУ им. Н.И. Лобачевского

\end{thebibliography}
\newpage

\appendix
\section{Приложение. Код}
\begin{lstlisting}[language=C++, frame=single]
#include "common/memory_utils.hpp"
#include "common/numeric_utils.hpp"
#include "common/test_utils.hpp"

#include <chrono>
#include <cstring>
#include <immintrin.h>
#include <iomanip>
#include <iostream>
#include <omp.h>
#include <random>
#include <stdlib.h>
#include <vector>

static constexpr int64_t maskarr[] = { -1, -1, -1, -1, 0, 0, 0, 0 };

struct SubMatrix {
    double* ptr;
    int n;
    int m;
    int stride;
};

class MatBuffer {
public:
    explicit MatBuffer(size_t size, size_t alignment = 32) : _size(numeric::divUp(size * sizeof(double), alignment) * alignment), _align(alignment) {
        _buffer = memory::aligned_alloc(_size, _align);
    }
    ~MatBuffer() {
        memory::aligned_free(_buffer);
    }

    MatBuffer(const MatBuffer&) = delete;
    MatBuffer& operator=(const MatBuffer&) = delete;

    double* data() {
        return static_cast<double*>(_buffer);
    }

    void memset(double val) {
        ::memset(_buffer, val, _size);
    }

private:
    void* _buffer;
    size_t _align;
    size_t _size;
};

//
// Main funcs
//

// Solve AXT = BT for X
inline void solveLinearSystemTransposed(const SubMatrix& A, const SubMatrix& B, SubMatrix& X) {
    const int packSize = 4;
#pragma omp parallel for
    for (int k = 0; k < B.n; ++k) {
        for (int j = 0; j < A.m; ++j) {
            double acc = B.ptr[j + B.stride * k];

            int quot = j / packSize;
            int rem = j % packSize;
            int X_start = X.stride * k;
            int A_start = A.stride * j;
            int pos = 0;
            __m256d sum = _mm256_setzero_pd();
            for (int s = 0; s < quot; ++s) {
                pos = s * packSize;
                __m256d A_elems = _mm256_loadu_pd(&A.ptr[pos + A_start]);
                __m256d X_elems = _mm256_loadu_pd(&X.ptr[pos + X_start]);
                sum = _mm256_fmadd_pd(A_elems, X_elems, sum);
            }
            pos = j - rem;

            __m256i mask = _mm256_loadu_si256((const __m256i*)(maskarr + (packSize - rem)));
            __m256d A_tail = _mm256_maskload_pd(&A.ptr[pos + A_start], mask);
            __m256d X_tail = _mm256_maskload_pd(&X.ptr[pos + X_start], mask);
            sum = _mm256_fmadd_pd(A_tail, X_tail, sum);

            sum = _mm256_hadd_pd(sum, sum);
            __m128d hi = _mm256_extractf128_pd(sum, 1);
            __m128d res = _mm_add_sd(_mm256_castpd256_pd128(sum), hi);
            acc -= _mm_cvtsd_f64(res);

            X.ptr[j + X.stride * k] = acc / A.ptr[j + A.stride * j];
        }
    }
}

inline void choleskyDecomp(const SubMatrix& A, SubMatrix& L) {
    constexpr int packSize = 4;
    for (int j = 0; j < A.m; ++j) {
        int L_start = L.stride * j;
        double acc = A.ptr[j + A.stride * j];
        int quot = j / packSize;
        int rem = j % packSize;

        __m256d sum = _mm256_setzero_pd();
        __m256d curr_val = _mm256_setzero_pd();
        for (int s = 0; s < quot; ++s) {
            curr_val = _mm256_loadu_pd(&L.ptr[s * packSize + L_start]);
            sum = _mm256_fmadd_pd(curr_val, curr_val, sum);
        }
        __m256i mask = _mm256_loadu_si256((const __m256i*)(maskarr + (packSize - rem)));
        curr_val = _mm256_maskload_pd(&L.ptr[quot * packSize + L_start], mask);
        sum = _mm256_fmadd_pd(curr_val, curr_val, sum);

        sum = _mm256_hadd_pd(sum, sum);
        __m128d hi = _mm256_extractf128_pd(sum, 1);
        __m128d res = _mm_add_sd(_mm256_castpd256_pd128(sum), hi);
        acc -= _mm_cvtsd_f64(res);

        acc = std::sqrt(acc);
        L.ptr[j + L_start] = acc;

#pragma omp parallel for private(acc, quot, rem, sum) schedule(dynamic, 8)
        for (int s = j + 1; s < L.n; ++s) {
            acc = A.ptr[j + A.stride * s];

            quot = j / packSize;
            rem = j % packSize;
            sum = _mm256_setzero_pd();
            for (int k = 0; k < quot; ++k) {
                __m256d a = _mm256_loadu_pd(&L.ptr[k * packSize + L_start]);
                __m256d b = _mm256_loadu_pd(&L.ptr[k * packSize + L.stride * s]);
                sum = _mm256_fmadd_pd(a, b, sum);
            }
            __m256i mask = _mm256_loadu_si256((const __m256i*)(maskarr + (packSize - rem)));
            __m256d a_tail = _mm256_maskload_pd(&L.ptr[quot * packSize + L_start], mask);
            __m256d b_tail = _mm256_maskload_pd(&L.ptr[quot * packSize + L.stride * s], mask);
            sum = _mm256_fmadd_pd(a_tail, b_tail, sum);

            sum = _mm256_hadd_pd(sum, sum);
            __m128d hi = _mm256_extractf128_pd(sum, 1);
            __m128d res = _mm_add_sd(_mm256_castpd256_pd128(sum), hi);
            acc -= _mm_cvtsd_f64(res);

            L.ptr[j + L.stride * s] = acc / L.ptr[j + L_start];
        }
    }
}

inline void extractBlock(const SubMatrix& block, SubMatrix& res) {
    for (int i = 0; i < res.n; ++i) {
        memcpy(&res.ptr[res.stride * i], &block.ptr[block.stride * i], res.m * sizeof(decltype(*block.ptr)));
    }
}

inline void storeBlock(SubMatrix& block, const SubMatrix& res) {
    for (int i = 0; i < res.n; ++i) {
        memcpy(&block.ptr[block.stride * i], &res.ptr[res.stride * i], res.m * sizeof(decltype(*block.ptr)));
    }
}

inline void storeBlockTransposed(SubMatrix& block, const SubMatrix& res) {
#pragma omp parallel for
    for (int i = 0; i < res.m; ++i) {
        for (int j = 0; j < res.n; ++j) {
            block.ptr[j + block.stride * i] = res.ptr[i + res.stride * j];
        }
    }
}

// X = ABT
inline void transposedMatMulRef(const SubMatrix& A, const SubMatrix& B, SubMatrix& X) {
#pragma omp parallel for
    for (int i = 0; i < A.n; ++i) {
        for (int j = 0; j < B.n; ++j) {
            for (int k = 0; k < A.m; ++k) {
                X.ptr[j + X.stride * i] += A.ptr[k + A.stride * i] * B.ptr[k + B.stride * j];
            }
        }
    }
}

inline void transposedMatMul(const SubMatrix& A, const SubMatrix& B, SubMatrix& X) {
    constexpr int pack_size = 4;
#pragma omp parallel for
    for (int i = 0; i < A.n; ++i) {
        auto X_start = X.stride * i;
        auto A_start = A.stride * i;
        for (int j = 0; j < B.n; ++j) {
            auto B_start = B.stride * j;

            __m256d sum = _mm256_setzero_pd();
            __m256d a, b;
            int idx = 0;
            for (; idx <= A.m - pack_size; idx += pack_size) {
                a = _mm256_loadu_pd(&A.ptr[idx + A_start]);
                b = _mm256_loadu_pd(&B.ptr[idx + B_start]);
                sum = _mm256_fmadd_pd(a, b, sum);
            }
            __m256i mask = _mm256_loadu_si256((const __m256i*)(maskarr + (pack_size - (A.m - idx))));
            __m256d a_tail = _mm256_maskload_pd(&A.ptr[idx + A_start], mask);
            __m256d b_tail = _mm256_maskload_pd(&B.ptr[idx + B_start], mask);
            sum = _mm256_fmadd_pd(a_tail, b_tail, sum);
            sum = _mm256_hadd_pd(sum, sum);
            __m128d hi = _mm256_extractf128_pd(sum, 1);
            __m128d res = _mm_add_sd(_mm256_castpd256_pd128(sum), hi);
            X.ptr[j + X_start] = _mm_cvtsd_f64(res);
        }
    }
}

// X = A - B
void matAdd(const SubMatrix& A, const SubMatrix& B, SubMatrix& X) {
    const int pack_size = 4;
    const int quot = A.m / pack_size;
    const int rem = A.m % pack_size;

#pragma omp parallel for
    for (int i = 0; i < A.n; ++i) {
        int pos = 0;
        int A_start = A.stride * i;
        int B_start = B.stride * i;
        int X_start = X.stride * i;

        __m256d res = _mm256_setzero_pd();
        for (int j = 0; j < quot; ++j) {
            pos = j * pack_size;
            __m256d a = _mm256_loadu_pd(&A.ptr[pos + A_start]);
            __m256d b = _mm256_loadu_pd(&B.ptr[pos + B_start]);
            res = _mm256_add_pd(a, b);
            _mm256_storeu_pd(&X.ptr[pos + X_start], res);
        }
        pos = quot * pack_size;
        __m256i mask = _mm256_loadu_si256((const __m256i*)(maskarr + (pack_size - rem)));
        __m256d a_tail = _mm256_maskload_pd(&A.ptr[pos + A_start], mask);
        __m256d b_tail = _mm256_maskload_pd(&B.ptr[pos + B_start], mask);
        res = _mm256_add_pd(a_tail, b_tail);
        _mm256_maskstore_pd(&X.ptr[pos + X_start], mask, res);
    }
}

// X = A - B
void matSub(const SubMatrix& A, const SubMatrix& B, SubMatrix& X) {
    const int pack_size = 4;
    const int quot = A.m / pack_size;
    const int rem = A.m % pack_size;

#pragma omp parallel for
    for (int i = 0; i < A.n; ++i) {
        int pos = 0;
        int A_start = A.stride * i;
        int B_start = B.stride * i;
        int X_start = X.stride * i;

        __m256d res = _mm256_setzero_pd();
        for (int j = 0; j < quot; ++j) {
            pos = j * pack_size;
            __m256d a = _mm256_loadu_pd(&A.ptr[pos + A_start]);
            __m256d b = _mm256_loadu_pd(&B.ptr[pos + B_start]);
            res = _mm256_sub_pd(a, b);
            _mm256_storeu_pd(&X.ptr[pos + X_start], res);
        }
        pos = quot * pack_size;
        __m256i mask = _mm256_loadu_si256((const __m256i*)(maskarr + (pack_size - rem)));
        __m256d a_tail = _mm256_maskload_pd(&A.ptr[pos + A_start], mask);
        __m256d b_tail = _mm256_maskload_pd(&B.ptr[pos + B_start], mask);
        res = _mm256_sub_pd(a_tail, b_tail);
        _mm256_maskstore_pd(&X.ptr[pos + X_start], mask, res);
    }
}

void blockedCholeskyDec(const SubMatrix& A, SubMatrix& L) {
    const int blk_mult = 512;
    const int blk_size = A.n < blk_mult ? A.n : blk_mult;
    const int blk_buffer_size = blk_size * blk_size;
    const int rem_blk_size = A.n - blk_size;
    const int rem_ll_buffer_size = rem_blk_size * blk_size;
    const int rem_lr_buffer_size = rem_blk_size * rem_blk_size;
    MatBuffer A_ul_buf(blk_buffer_size);
    SubMatrix A_ul{ A_ul_buf.data(), blk_size, blk_size, blk_size };
    SubMatrix A_ul_src{ A.ptr, blk_size, blk_size, A.stride };
    MatBuffer L_ul_buf(blk_buffer_size);
    SubMatrix L_ul{ L_ul_buf.data(), blk_size, blk_size, blk_size };

    extractBlock(A_ul_src, A_ul);
    choleskyDecomp(A_ul, L_ul);

    SubMatrix L_ul_dst{ L.ptr, blk_size, blk_size, L.stride };
    storeBlock(L_ul_dst, L_ul);

    if (rem_blk_size == 0) {
        return;
    }

    MatBuffer A_ll_buf(rem_ll_buffer_size);
    SubMatrix A_ll{ A_ll_buf.data(), rem_blk_size, blk_size, blk_size };
    SubMatrix A_ll_src{ &A.ptr[blk_size * A.stride], rem_blk_size, blk_size, A.stride };
    MatBuffer L_ll_buf(rem_ll_buffer_size);
    SubMatrix L_ll{ L_ll_buf.data(), rem_blk_size, blk_size, blk_size };

    extractBlock(A_ll_src, A_ll);
    solveLinearSystemTransposed(L_ul, A_ll, L_ll);
    SubMatrix L_ll_dst{ &L.ptr[blk_size * L.stride], rem_blk_size, blk_size, L.stride };
    storeBlock(L_ll_dst, L_ll);

    MatBuffer A_lr_buf(rem_lr_buffer_size);
    SubMatrix A_lr{ A_lr_buf.data(), rem_blk_size, rem_blk_size, rem_blk_size };
    SubMatrix A_lr_src{ &A.ptr[blk_size * A.stride + blk_size], rem_blk_size, rem_blk_size, A.stride };

    extractBlock(A_lr_src, A_lr);

    MatBuffer L_ll_prod_buf(rem_lr_buffer_size);
    SubMatrix L_ll_prod{ L_ll_prod_buf.data(), rem_blk_size, rem_blk_size, rem_blk_size };

    transposedMatMul(L_ll, L_ll, L_ll_prod);
    matSub(A_lr, L_ll_prod, A_lr);

    SubMatrix L_lr{ &L.ptr[blk_size * L.stride + blk_size], rem_blk_size, rem_blk_size, L.stride };
    blockedCholeskyDec(A_lr, L_lr);
}

void Cholesky_Decomposition(double* A, double* L, int n) {
    SubMatrix sA{ A, n, n, n };
    SubMatrix sL{ L, n, n, n };

    blockedCholeskyDec(sA, sL);
}

//
// Helpers
//

void randomlyFillLowerTriang(SubMatrix& X) {
    std::random_device rd;
    std::mt19937 gen(time(0));
    double max = 1.5l, min = 1.l;
    std::uniform_int_distribution<int> dist(0, 100);

    for (size_t i = 0; i < X.m; i++) {
        for (size_t j = 0; j < i + 1; j++) {
            X.ptr[j + X.stride * i] = static_cast<double>(dist(gen)) / 100.l * (max - min) + min;
        }
    }
}

double calRelResidual(const SubMatrix& A, const SubMatrix& X) {
    // Cacl relative residual using Frobenius norm
    double acc = 0.l;
    double t;
    for (size_t i = 0; i < A.m; ++i) {
        for (size_t j = 0; j < A.n; ++j) {
            t = A.ptr[j + A.stride * i] - X.ptr[j + X.stride * i];
            if (std::isnan(t)) {
                t = 0.l;
            }
            acc += t * t;
        }
    }

    auto diff_norm = std::sqrt(acc);

    acc = 0.l;
    for (size_t i = 0; i < A.m; ++i) {
        for (size_t j = 0; j < A.n; ++j) {
            t = A.ptr[j + A.stride * i];
            acc += t * t;
        }
    }

    auto A_norm = std::sqrt(acc / A.m);

    if (A_norm == 0) {
        return 0;
    }

    return diff_norm / A_norm;
}
\end{lstlisting}
\end{document}